\chapter{\texorpdfstring{$A_\infty$}{A∞} 范畴}
\label{chapter:I}
\SourcePage{15}
本章介绍 Fukaya 范畴背后的同调代数。我们先快速回顾 $\Ai$ 结构。这绝不是对该理论的完整阐述。基本上，唯一详细讨论的方面是恒等态射，因为它对 Fukaya 范畴而言可能颇为棘手。此外，我们有意使形式体系保持具体。这意味着我们不使用操作子（operad）或非交换几何，尽管这两种语言中的任何一种，都可以帮助理清错综复杂的公式。\znote{7}{操作子（operad）是用\textquotedblleft $n$ 元运算的合成法则\textquotedblright 编码代数结构的语言：$\Ai$ 操作子是结合操作子的一个 dg 分辨，其上的代数恰为 $\Ai$ 代数，$\mu^d$ 就是 $d$ 元运算。本书不用这套语言，但如\stref{命题}{1.12}那样的递归构造，用操作子说就是\textquotedblleft 把分辨搬到范畴上\textquotedblright，公式里的树形求和因此有了统一解释。[来源：nLab \textquotedblleft A-infinity-algebra\textquotedblright，https:/\allowbreak /\allowbreak ncatlab.org/\allowbreak nlab/\allowbreak show/\allowbreak A-infinity-algebra（2026-11-15 访问）。]}关于这些材料，还有风格各异的其他参考文献。这里给出几点建议：一些较早的基础论文，例如\bcite{70}、\bcite{62}，即使重点与本书稍有不同，仍然值得一读。在较新的文献中，Keller 的工作\bcite{82}、\bcite{83}、\bcite{84}既对基础材料作了易于理解的说明，也讨论了它们在代数和表示论问题中的应用。Lefèvre 的学位论文\bcite{95}采用同伦代数语言，从相近的观点作了全面阐述。Fukaya 及其合作者写过多份专门面向辛几何的论述\bcite{51}、\bcite{52}、\bcite{53}、\bcite{58}、\bcite{60}、\bcite{56}、\bcite{55}。最后三篇文献的一般性远远超过我们的需要，但愿意顾及这一点的读者会发现，许多论证可以简化，适用于我们较为有限的框架。最后，Kontsevich 的论文和讲义\bcite{90}、\bcite{88}、\bcite{89}中也蕴藏着丰富的见解。

本章其余部分的总主题，是 $\Ai$ 范畴中对象的构造与运用。我们将讨论映射锥和正合三角，以及将对象分裂成直和项的问题。还会考察一种更具体的构造，即所谓一个对象沿另一个对象的扭转。此前，这一构造在文献中出现于两个不同背景。其一是例外集与变异的概念；它最初是为在代数几何中构造刚性向量丛而产生的：参看\bcite{67}和论文集\bcite{117}。这一理论逐步推广到一般三角范畴（不同方面见\bcite{65}、\bcite{26}、\bcite{27}、\bcite{28}）；文献\bcite{66}是一篇有用的新近综述。扭转构造第二次、也是较近的一次出现，与球面对象有关；在这一背景下，它给出导出范畴的自同构\bcite{134}。向 $\Ai$ 范畴的推广由 Kontsevich 开创，我们将广泛借鉴他的思想。除了建立一般形式体系外，我们还将推出若干结论，特别是\stref{推论}{5.8}和\stref{命题}{5.17}；后面把这些材料与辛几何及 Lefschetz 纤维化联系起来时，它们会派上用场。本章最后以一段略带技术性的讨论收尾，简要介绍具有 $\Z/2$ 对称性的 $\Ai$ 范畴。
